% ISOLATED FUTURE INTEGRATION CANDIDATE. See GATE.md; not a live proof replacement.

\subsection{The skeleton, its lengths, and the Laplacian}
\label{sec:ch31-curve-setup}

In this subsection and through \cref{rem:ch31-curve-solution-scope},
\(K\) is any complete nontrivially valued non-Archimedean field.
Let \(X\) be a smooth projective geometrically connected \(K\)-curve,
and suppose a proper strictly semistable formal model \(\cX\) over
\(K^\circ\) is given.  Write \(\Gamma\) for its finite connected
skeleton and \(p_{\cX}:X^{\An}\to\Gamma\) for its retraction.
No assertion that every curve has such a model over the original
\(K\) is implicit in this hypothesis.

An edge \(e\), corresponding to a node \(p_e\) of the special fibre,
has its analytic modulus length \(\ell_e>0\) and conductance
\[
                    c_e=[\kappa(p_e):\widetilde K].
\]
The conductance is the same at both ends, including both incidences
of a loop, and is retained under subdivision
\cite[\S2.2.3, p.~41]{Thu05}.
Lengths are measured by the chosen absolute value, using the natural
logarithm: an annulus \(r<|t|<R\) has length \(\log(R/r)\).
Where an \'{e}tale node chart has equation \(xy=a\), this is
\(-\log|a|\).  A node need not be rational over \(\widetilde K\);
one must use its \'{e}tale chart, not silently replace its completed
local ring by a rational-node equation.  In the discretely valued
examples with \(-\log|\pi|=1\), \(xy=\pi^a\) gives length \(a\).
Regularity of that total-space node forces \(a=1\); general
semistable formal models need not be regular before resolution.

For a real continuous piecewise affine function \(F\) on \(\Gamma\),
define
\begin{equation}
\label{eq:ch31-graph-laplacian}
 \Delta_\Gamma F
   :=\sum_{x\in\Gamma}
       \left(\sum_{\nu\in T_x\Gamma}
              c_{e(\nu)}\partial_\nu F(x)\right)\delta_x,
 \qquad
 \mathrm d\mathrm d^c(F\circ p_{\cX})=\iota_*\Delta_\Gamma F ,
\end{equation}
where \(\iota:\Gamma\hookrightarrow X^{\An}\), and
\(\partial_\nu\) is the outgoing derivative in the analytic length
coordinate.  The second equality is
\cite[Theorem~3.2.10(ii)]{Thu05}; properness removes its boundary
exception.  It also says that \(F\circ p_{\cX}\) is Thuillier-lisse.
It does not say that this function is a rational model function.
That stronger assertion requires the slope, breakpoint, and value-group
compatibilities of \cite[Proposition~3.5]{GJKM19}.

Edgewise integration by parts, counting each edge once, gives
\begin{equation}
\label{eq:graph-dirichlet-sign}
 \Delta_\Gamma F(\Gamma)=0,\qquad
 \int_\Gamma F\,\mathrm d\Delta_\Gamma F
       =-\sum_e c_e\int_e|F'|^2\,\mathrm ds\leq0.
\end{equation}
The two endpoint terms of every edge cancel in the total mass.
For the second identity they combine with the interior slope jumps
to give the displayed integral; equivalently use the bilinear
identity proved in \cref{lem:ch31-green-kernel}.
The weights are essential even when \(X\) is geometrically connected.

\begin{remark}[Lengths, subdivisions, and attached branches]
\label{rem:ch31-lengths}
The combinatorial graph depends on the model.  An admissible
semistable refinement can subdivide edges and can also attach new
branches.  More precisely, the old skeleton embeds in the new one,
their retractions commute, each new connected piece attaches at one
point, and the weighted Laplacian pushes forward under the induced
retraction \cite[Proposition~2.2.26]{Thu05}.
Subdivision leaves both the analytic edge length and its conductance
unchanged on the subdivided pieces.  The resistance length
\(\ell_e/c_e\) used for the estimates below is an auxiliary metric;
it does not replace the analytic modulus length in the slope formula.

A model function determined on a dominating model can vary on an
attached branch.  It need not factor through the original \(p_{\cX}\).
One must either impose that factorization explicitly or compute on
a skeleton that actually determines the function and the reference.
A common semistable model can be used after the required scalar
extension; its existence over the original \(K\) is not presumed.
\end{remark}

\subsection{The Monge--Amp\`ere operator is affine}
\label{sec:ch31-curve-affine}

Let \(L\) have degree \(V>0\).  In this subsection assume an actual
relatively nef formal line bundle \(\cL\) on the given \(\cX\)
extends \(L^{\otimes m}\), for an integer \(m>0\).
Let \(\lVert\cdot\rVert_0\) be its rational formal metric and
\(\omega=c_1(L,\lVert\cdot\rVert_0)\).
This is additional data, not a consequence of the existence of
\(\cX\).  Put \(\mu_0=\MA_\omega(0)\).

\begin{lemma}
\label{lem:ch31-curve-ma-affine}
The reference measure is
\[
 \mu_0=\sum_v d_v\delta_v,\qquad
 d_v=\frac1m\deg_{\widetilde K}(\cL|_{C_v}),\qquad
 \sum_vd_v=V.
\]
It is positive and supported on the original vertex set.
For every rational model function of the form
\(u=F\circ p_{\cX}\), where \(F\) is real piecewise affine,
\begin{equation}
\label{eq:curve-ma-laplacian}
                 \MA_\omega(u)=\mu_0+\iota_*\Delta_\Gamma F.
\end{equation}
Here for a model metric that is not nef the left side means its signed
curvature measure.  The function \(u\) is \(\omega\)-psh in the
model sense if and only if \(\mu_0+\Delta_\Gamma F\geq0\).
For a measure \(\mu\) supported on \(\Gamma\), the equation becomes
\begin{equation}
\label{eq:curve-poisson}
                      \Delta_\Gamma F=\mu-\mu_0 .
\end{equation}
Measures on \(\Gamma\) are identified with their pushforwards on
\(X^{\An}\) when no confusion is possible.

More generally, let \(\theta\) be a real closed model form determined
on \(\cX\), let \(d_x(\theta)\) be its component degree over
\(\widetilde K\) at an original vertex, and put \(d_x(\theta)=0\)
elsewhere.  For an actual rational model function
\(F\circ p_{\cX}\), the numerical criterion is
\begin{equation}
\label{eq:curve-slope-criterion}
 \theta+\mathrm d\mathrm d^c(F\circ p_{\cX})\text{ is nef}
 \quad\Longleftrightarrow\quad
 d_x(\theta)+\sum_{\nu\in T_x\Gamma}
       c_{e(\nu)}\partial_\nu F(x)\geq0
 \quad\text{for every }x\in\Gamma.
\end{equation}
Only the original vertices and slope breaks impose nonzero conditions.
\end{lemma}

\begin{proof}
For the reference model the curvature mass at \(v\) is its degree
over the base residue field divided by \(m\).
In \cite[Theorem~18.5.5 and Proposition~18.5.6(a)--(c)]{CLD25},
the generic smoothness and geometrically reduced strictly semistable
special fibre make the additional multiplicity, defect, and graded
residue factors equal to one.  Total mass is \(V\) by
\cite[Theorem~18.3.11]{CLD25}.
If a component has constant field \(k_v\), its degree over
\(\widetilde K\) is \([k_v:\widetilde K]\) times its degree over
\(k_v\); the latter cannot be substituted without this factor.

The formal weights are Thuillier-lisse by the tensor-power and
analytic-unit argument in the proof of
\cref{thm:curves-arbitrary-field-continuity}.
The original-field model comparison there identifies Thuillier and
CLD curvature, so additivity and \eqref{eq:ch31-graph-laplacian}
prove \eqref{eq:curve-ma-laplacian}.
For the positivity criterion, a nef model metric has positive
curvature.  Conversely, positive curvature makes each lisse local
weight Thuillier-subharmonic by \cite[Theorem~3.4.12(ii)]{Thu05}.
This remains true after extension to \(\widehat{K^{\mathrm a}}\)
by \cite[Corollary~3.4.5]{Thu05}.  There one may choose a common
semistable model actually determining the pulled-back metric; the
slope and component-degree formula
\cite[Proposition~3.7]{GJKM19} says all its component degrees are
nonnegative.  Thus that model metric is nef.
Model-independence and reflection of nefness under scalar extension
\cite[Proposition~3.5 and Lemma~3.3]{GM19} give nefness over \(K\).
Finally \eqref{eq:curve-poisson} is
\eqref{eq:curve-ma-laplacian} rearranged.

For the real-form criterion \eqref{eq:curve-slope-criterion},
the component-degree and outgoing-slope identities extend real-linearly
from model line bundles.  On a common semistable model over
\(\widehat{K^{\mathrm a}}\) determining the finitely many classes
representing \(\theta\) and the actual function \(F\circ p_{\cX}\),
nefness means that every special-fibre component degree is nonnegative.
This is the real-form criterion of \cite[Proposition~3.7]{GJKM19}.
Its original-field normalization uses the residue-degree weights
in \eqref{eq:ch31-graph-laplacian}.  Nefness reflects over \(K\)
for real numerical classes by the degree and projection-formula
argument in the proof of
\cref{thm:ch29-real-class-curve-model-envelope}.
Thus the same criterion holds over \(K\).
This argument never extends \(F\) by an assertion of factorization
for all functions on dominating models.
\end{proof}

\begin{warning}{Affine, not linear}
Equation \eqref{eq:curve-ma-laplacian} is affine in the potential,
with linear part \(\Delta_\Gamma\) and constant term \(\mu_0\).
The set of continuous semipositive metric weights relative to
\(\lVert\cdot\rVert_0\) is a convex set, not in general a cone or a
vector space: the constraint is
\(\mu_0+\Delta_\Gamma F\geq0\).
Within the rational model class, taking an arbitrary real convex
combination need not remain a rational model function.
\end{warning}

\subsection{Solving the equation: the graph Green kernel}
\label{sec:ch31-curve-green}

The following graph lemma uses arbitrary positive conductances and
does not depend on any analytic existence or envelope theorem.

\begin{lemma}
\label{lem:ch31-green-kernel}
Let \(\Gamma\) be a finite connected metric graph.  Give each edge \(e\)
a positive real constant \(c_e\), retained on every subdivision of that
edge.  Use the weighted outgoing-slope Laplacian
\eqref{eq:ch31-graph-laplacian}.
Let \(q\in \Gamma\).  There is a unique symmetric continuous kernel
\(g_q:\Gamma\times \Gamma\to\mathbf R\) for which
\begin{equation}
\label{eq:ch31-green-kernel}
 \Delta_{\Gamma,x}g_q(x,y)=\delta_y-\delta_q,
 \qquad g_q(q,y)=0,
\end{equation}
and \(g_q(\cdot,y)\) is affine outside the graph vertices and
\(\{q,y\}\).  If \(d_R\) is the path metric assigning length
\(\ell_e/c_e\) to \(e\), and \(D_R=\operatorname{diam}_{d_R}\Gamma\), then
\begin{equation}
\label{eq:ch31-green-kernel-estimates}
 |g_q(x,y)|\leq D_R,
 \qquad
 |g_q(x,y)-g_q(x',y)|\leq\sqrt{D_Rd_R(x,x')}.
\end{equation}
For a finite signed Radon measure \(\nu\) of mass zero,
\begin{equation}
\label{eq:ch31-green-potential}
 U_\nu(x)=\int_\Gamma g_q(x,y)\,\mathrm d\nu(y)
\end{equation}
is the unique continuous solution of \(\Delta_\Gamma U_\nu=\nu\) with
\(U_\nu(q)=0\).  Distributional equality here means
\(\int_\Gamma U_\nu\,\mathrm d\Delta_\Gamma F=\int_\Gamma F\,\mathrm d\nu\)
for every continuous piecewise affine \(F\).
If \(\nu_j\to\nu\) weakly and their total variations are uniformly
bounded, then \(U_{\nu_j}\to U_\nu\) uniformly.
\end{lemma}


\begin{proof}
Subdivide to include all breakpoints.  Edgewise integration by parts gives
\begin{equation}
\label{eq:ch31-weighted-green-energy}
 \int_\Gamma F\,\mathrm d\Delta_\Gamma H
   =-\sum_e c_e\int_e F'H'\,\mathrm ds
   =\int_\Gamma H\,\mathrm d\Delta_\Gamma F.
\end{equation}
In particular \(\Delta_\Gamma F(\Gamma)=0\).  An edge is counted once in the
energy and both of its tangent directions are counted in the vertex
sum, including for loops.  This is the weighted version of
\cite[Proposition~1.2.13]{Thu05}; it also follows immediately by
cancelling the two endpoint terms of each edge.

Fix \(y\), and insert \(q,y\) as vertices.  On functions affine on
each resulting edge, minus the Laplacian has quadratic form
\[
                 \sum_{e=(v,w)}\frac{c_e}{\ell_e}
                                      (a_v-a_w)^2.
\]
The associated matrix is symmetric; its kernel consists exactly of
constants by connectedness.  Its image is consequently the vectors
whose coordinate sum is zero.  Solve for \(\delta_y-\delta_q\),
and normalize the solution to vanish at \(q\).  Further subdivisions
do not change it, since the affine extension has zero mass at each
new point.  If \(\Gamma\) is a point, the solution is identically zero.

Suppose now that \(q\ne y\).  The maximum principle for this finite
weighted matrix says that the maximum of \(h=g_q(\cdot,y)\) occurs
at \(q\) and its minimum at \(y\).  Indeed at a maximum every
outgoing slope is nonpositive.  If the Laplacian there is zero,
positivity of every conductance forces all those slopes to vanish;
propagate the maximum along edges until meeting a nonzero mass.
A positive mass cannot occur at a maximum.  The argument for a
minimum is the same with signs reversed.  Thus \(h(q)=0\),
\(h\leq0\), and \eqref{eq:ch31-weighted-green-energy} gives
\[
 E(h):=\sum_e c_e\int_e|h'|^2\,\mathrm ds=-h(y)
                                         =\operatorname{osc}(h).
\]
For any edge-simple path \(\gamma\) from \(a\) to \(b\), after
subdivision at its endpoints, Cauchy--Schwarz gives
\[
 |h(a)-h(b)|^2
 \leq\left(\int_\gamma\frac{\mathrm ds}{c_e}\right)
       \left(\int_\gamma c_e|h'|^2\,\mathrm ds\right)
 \leq\operatorname{length}_R(\gamma)E(h).
\]
The second inequality uses that no edge segment is traversed twice.
Shortest resistance-length paths may be chosen edge-simple, since all
edge lengths are positive.  Take such a path from \(q\) to \(y\) to
obtain \(\operatorname{osc}(h)\leq d_R(q,y)\leq D_R\), then use
a shortest path from \(x\) to \(x'\).  This proves
\eqref{eq:ch31-green-kernel-estimates}.

Apply \eqref{eq:ch31-weighted-green-energy} to
\(g_q(\cdot,y)\) and \(g_q(\cdot,z)\), on a subdivision containing
both \(y,z\).  The normalizing values at \(q\) vanish, giving
\(g_q(z,y)=g_q(y,z)\).  Symmetry and the uniform estimate already
proved give the same modulus of continuity in the second variable.
Hence the kernel is jointly continuous.  This avoids assuming
continuous dependence of a moving-vertex matrix at a degenerating edge.

The integral in \eqref{eq:ch31-green-potential} is therefore
continuous.  For piecewise affine \(F\), Fubini and
\eqref{eq:ch31-weighted-green-energy} yield
\[
 \int_\Gamma U_\nu\,\mathrm d\Delta_\Gamma F
  =\int_\Gamma\left(F(y)-F(q)\right)\mathrm d\nu(y)
  =\int_\Gamma F\,\mathrm d\nu.
\]
Both integrations are against finite signed measures and a bounded
continuous kernel, so Fubini applies.

For uniqueness in the asserted \emph{continuous} class, do not assume
an energy integral before proving piecewise affineness.  If \(w\) is
continuous and distributionally harmonic, test with the piecewise
affine tent on a closed subinterval \([a,b]\) of an open edge, of
height one at \(x\in(a,b)\) and zero outside \((a,b)\).
The equality \(\int w\,\mathrm d\Delta_\Gamma F=0\) becomes
\[
 \frac{w(a)}{x-a}+\frac{w(b)}{b-x}
       -w(x)\left(\frac1{x-a}+\frac1{b-x}\right)=0.
\]
It follows that \(w\) is affine on every open edge, and continuity
extends this to its endpoints.  Vertex tent tests then give the
weighted zero-slope sum at the vertices.  The energy identity now
applies and forces \(w\) to be constant.

Finally, the family \(\{g_q(x,\cdot):x\in \Gamma\}\) is compact in
\(C(\Gamma)\), being the continuous image of the compact graph.  Weak
convergence of measures with a uniform total-variation bound is uniform
on this compact family: approximate it by a finite sup-norm net and
bound the error by that total-variation bound.  This proves the last
assertion.
\end{proof}



\begin{proposition}
\label{prop:graph-ma-solution}
With the given model reference above, let \(\mu\) be any positive Radon
measure on \(\Gamma\) of mass \(V\).  The function
\(F=U_{\mu-\mu_0}\) of \cref{lem:ch31-green-kernel} is the unique
continuous solution of \eqref{eq:curve-poisson} normalized by \(F(q)=0\).
The metric
\[
             \lVert\cdot\rVert_0\mathrm e^{-F\circ p_{\cX}}
\]
is continuous and semipositive, and its Thuillier, Chambert--Loir,
and CLD curvature measures all equal \(\mu\).
The solution is unique up to an additive constant in the weight,
or a positive constant scale in the metric.
\end{proposition}

\begin{proof}
The measure \(\nu=\mu-\mu_0\) has mass zero.  The graph lemma gives
the asserted continuous potential and its distributional uniqueness.
To check its analytic curvature without assuming rational slopes,
choose finite positive atomic measures \(\mu_j\) of mass exactly
\(V\) converging weakly to \(\mu\).  Partition the finitely many
edges into intervals of mesh tending to zero, place each interval's
mass at one point, and retain vertex masses; this gives such measures
with real, not necessarily rational, masses.
Then \(F_j=U_{\mu_j-\mu_0}\) are real piecewise affine and converge
uniformly to \(F\) by \cref{lem:ch31-green-kernel}.

The reference weights and \(F_j\circ p_{\cX}\) are Thuillier-lisse.
The slope formula gives curvature \(\mu_j\) for the corresponding
metrics.  Their weights are subharmonic, and the uniform limits
remain subharmonic by the harmonic comparison definition.
For a compactly supported lisse test function \(a\), the error in
curvature pairing is bounded by
\(\lVert F_j-F\rVert_\infty
  \lVert\mathrm d\mathrm d^ca\rVert_{\mathrm{TV}}\).
Thus the limiting metric has Thuillier curvature \(\mu\);
density of lisse tests \cite[Proposition~3.2.7]{Thu05} identifies
the measures.

The optional model-reference formula
\eqref{eq:ch29-curve-model-reference-solution} in
\cref{thm:curves-arbitrary-field-continuity} now applies to exactly
this metric.  Its proof promotes the continuous positive-curvature
metric to a uniform limit of semipositive rational formal metrics
by comparison over the completed algebraic closure and finite
normal descent, and identifies the three curvature measures.
This use is not circular: that proof uses only the graph lemma above,
not the present proposition.
Uniqueness of analytic metrics follows because the logarithm of
their ratio is a continuous harmonic function on the proper connected
curve, hence constant, as proved in
\cref{prop:ch29-curve-point-reference}.

Neither the \(F_j\) nor their pullbacks have been declared rational
model functions.  Rational atomic masses alone would not imply
rational slopes on a graph with arbitrary modulus lengths.
For a nonatomic measure the limiting potential is only asserted
continuous; it need not be Thuillier-lisse.
\end{proof}

\begin{remark}[What this proposition is, and is not]
\label{rem:ch31-curve-solution-scope}
The proposition solves the equation for every positive measure of
the right mass supported on this given finite skeleton.
The quantitative bound is
\(\sup|g_q|\leq D_R\), where \(D_R\) is the diameter for resistance
lengths \(\ell_e/c_e\); the analytic lengths remain \(\ell_e\).
Measures not supported on a finite skeleton are outside this
continuous-solution assertion.
No degree-zero rational-model existence claim, unrestricted envelope
identity, or higher-dimensional reduction is inferred.
\end{remark}
